\documentclass[12pt]{article}

\usepackage{sbc-template}
\usepackage{graphicx,url}
\usepackage[utf8]{inputenc}
\usepackage[brazilian]{babel}
     
\sloppy

\title{Identificar e caracterizar perfis de municípios paulistas segundo seus padrões de criminalidade, relacionando-os a indicadores de efetividade da gestão municipal e de vulnerabilidade socioeconômica e urbana}

\author{Bruno Hideo Ioneda\inst{1}, Guilherme Samuel Lemos Segura\inst{1}, Higor Ranel Viani Lopes\inst{1}}


\address{Escola de Artes, Ciências e Humanidades -- Universidade de São Paulo
  (USP) \\ Rua Arlindo Bettio, 1000 -- 03828-000 -- São Paulo -- SP -- Brazil
  \email{brunoioneda@usp.br, guilhermesegura@usp.br,
  higor.ranel@usp.br}
}

\begin{document} 

\maketitle

\begin{abstract} 
The heterogeneity of crime and urban vulnerability in Brazil limits the effectiveness of standardized security policies. To address this challenge, this study aims to identify and characterize profiles of municipalities in the state of São Paulo, cross-referencing criminal patterns with indicators of municipal management effectiveness and socioeconomic and urban vulnerability. The methodology consolidates public data (IEGM/TCE-SP, SSP-SP, and IBGE) subjected to unsupervised clustering algorithms (K-means, Agglomerative Hierarchical, and DBSCAN), evaluated by internal validation metrics and a qualitative analysis.
\end{abstract}
   
\begin{resumo}
A heterogeneidade da criminalidade e da vulnerabilidade urbana
no Brasil limita a eficácia de políticas de segurança padronizadas. Para enfrentar esse desafio, este estudo visa identificar e caracterizar perfis de municípios paulistas, cruzando padrões criminais com indicadores de efetividade da gestão municipal e de vulnerabilidade socioeconômica e urbana. A metodologia consolida dados públicos (SSP-SP, IEGM/TCE-SP e IBGE) submetidos a algoritmos de clusterização não supervisionada (K-means, Hierárquico Aglomerativo e DBSCAN), avaliados por métricas de validação interna e uma análise qualitativa. 
\end{resumo}

\section{Introdução}
Estudos recentes mostram que a violência letal deixou de ser um flagelo exclusivamente urbano e passou a se interiorizar pelo território nacional \cite{soaresfilho2024}, concentrando-se principalmente em regiões com baixos índices de escolaridade e renda, além de fragilidade do poder público \cite{minamisava2009}.
O estado de São Paulo, com sua heterogeneidade espalhada por seus 645 municípios, torna-se um recorte relevante para investigar o fenômeno da violência. Órgãos públicos disponibilizam, de forma aberta, dados sobre criminalidade \cite{ssp-sp}, efetividade da gestão municipal \cite{tce-sp} e condições socioeconômicas \cite{ibge-sidra} que, quando analisados de forma conjunta, permitem identificar padrões entre crime, capacidade institucional e vulnerabilidade social.

Na literatura científica, métodos de clusterização são frequentemente utilizados para revelar padrões latentes por meio da interação de múltiplas variáveis urbanas e sociais. Diversas pesquisas ilustram a aderência dessa abordagem à análise da criminalidade: \cite{costasilva2022}, por exemplo, segmentaram municípios em Pernambuco com base em registros criminais para orientar políticas públicas. Internacionalmente, \cite{wang2013} identificaram dinâmicas espaciais de crimes atreladas a fatores socioeconômicos, enquanto \cite{gonzalez2026} empregaram clusterização e PCA para investigar a correlação entre homicídios e outras ocorrências policiais no Peru. Adicionalmente, o desenho de pesquisa deste trabalho dialoga com estudos recentes sobre aprendizado supervisionado no contexto urbano brasileiro \cite{silva2025}, de forma a incorporar estratégias de clusterização e avaliação sistemática da acurácia dos agrupamentos.

Apesar do avanço dessas abordagens, poucos trabalhos investigam indicadores de criminalidade, efetividade institucional e condições socioeconômicas somente  no contexto do estado de São Paulo. Essa lacuna complica a compreensão de como fatores estruturais, institucionais e sociais se combinam para explicar a distribuição desigual da violência dentro do território estadual.

Diante disso, este trabalho tem como objetivo identificar e caracterizar perfis
de municípios paulistas segundo seus padrões de criminalidade, relacionando-os
a indicadores de efetividade da gestão municipal e de vulnerabilidade socioeconômica
e urbana. Para isso, será construída uma base de dados que integra indicadores de
criminalidade, gestão pública e condições socioeconômicas dos municípios do estado,
sobre a qual diferentes algoritmos de clusterização serão aplicados e comparados
por meio de métricas de validação interna e uma análise qualitativa. A caracterização dos perfis resultantes, fundamentada na literatura sobre desorganização social e vulnerabilidade urbana,
permitirá compreender como a violência se distribui e se relaciona com o contexto
institucional e socioeconômico dos municípios paulistas.


\section{Trabalhos relacionados}
Nesta análise, foram analisados quatro trabalhos dos quais se valeram de abordagens metodológicas próximas à proposta deste projeto, seja por meio de aprendizado não supervisionado imposto sobre dados de criminalidade, seja por meio de estatística espacial.

Soares Filho et al. \cite{soaresfilho2024} realizaram um estudo ecológico com a totalidade dos municípios brasileiros (2000--2018), de forma que relacionaram taxas de homicídio a indicadores de violência, sociodemográficos e de segurança pública. Utilizaram índices de Moran e regressão espacial (\textit{spatial lag}), os autores identificaram a taxa de eventos violentos, a proporção de população negra, o baixo nível educacional e a presença de estruturas policiais municipais como principais fatores associados aos homicídios, que evidenciou um deslocamento do fenômeno da região Sudeste para o Norte/Nordeste ao longo do período. Apesar da abrangência nacional, o estudo não emprega algoritmos de clusterização não supervisionada, tratando o homicídio como variável-resposta em modelos de regressão, e não incorpora indicadores de efetividade da gestão pública municipal.

Minamisava et al. \cite{minamisava2009} investigaram, na cidade de Goiânia, a distribuição espacial de mortes intencionais e não intencionais entre jovens de 15 a 24 anos, aplicando estatística de varredura espacial (\textit{spatial scan statistic}, via SaTScan) para identificar clusters de setores censitários com risco elevado. As características socioeconômicas do cluster mais provável foram comparadas ao restante do município, revelando forte associação entre desigualdade social e homicídios. O estudo é metodologicamente robusto e restrito a uma única cidade, porém não conta com a aplicação de técnicas de clusterização multivariada e com os indicadores que serão utilizados neste trabalho.

González de la Rosa \cite{gonzalez2026} analisaram a relação entre homicídios e 30 subtipos de ocorrências policiais nas províncias do Peru, combinando análise de correlação, clusterização hierárquica aglomerativa e Análise de Componentes Principais (PCA). Os autores identificaram três clusters criminais bem diferenciados, de forma a associar homicídios a crimes de motivação econômica. Essa é a abordagem mais próxima, em termos metodológicos, da proposta deste trabalho. Contudo, os autores utilizam apenas um algoritmo de clusterização, sem comparação com outras técnicas nem validação por métricas internas, além de não incorporarem indicadores de gestão pública ou de vulnerabilidade socioeconômica.

Por fim, Wang et al. \cite{wang2013} propuseram um arcabouço de mineração de dados espaciais baseado em \textit{Geospatial }Discriminative\textit{ Patterns} (GDPatterns) para identificar e otimizar hotspots de um crime-alvo (roubo residencial) a partir de variáveis relacionadas (taxa de prisão, densidade populacional, execuções hipotecárias, entre outras), utilizando agrupamento hierárquico aglomerativo para sumarizar os padrões descobertos em clusters. Aplicado a uma cidade norte-americana, o estudo concentra-se na otimização de fronteiras de hotspots para um único tipo de crime, e não incluiu a caracterização de perfis multivariados entre múltiplos municípios nem indicadores de efetividade municipal.

De forma geral, os trabalhos analisados evidenciam duas lacunas centrais que este projeto busca preencher. A priori, a maioria dos estudos concentra-se em estatística espacial (índices de Moran, varredura espacial) ou em um único algoritmo de clusterização, sem uma comparação sistemática entre diferentes técnicas (K-means, hierárquico aglomerativo e DBSCAN) avaliadas por múltiplas métricas de validação interna, como proposto neste trabalho. Em segundo lugar, nenhum dos estudos revisados integra, simultaneamente, indicadores de criminalidade, de efetividade da gestão municipal (IEGM/TCE-SP) e de vulnerabilidade socioeconômica (IBGE/SIDRA) em um recorte estadual específico. 

Este projeto busca, portanto, contribuir com uma caracterização mais completa dos perfis municipais de criminalidade no estado de São Paulo, de modo a visar a dimensão criminal à capacidade institucional e ao contexto socioeconômico dos municípios.


\section{Materiais e Métodos}
A metodologia do trabalho proposto segue em quatro etapas principais: aquisição de dados, pré processamento dos dados, aplicação e comparação de algoritmos de clusterização, e posteriormente caracterização dos perfis, tal  metodologia foi inspirada na estrutura adota por \cite{silva2025} (municípios referenciais, aquisição de dados, , e treinamento de modelo) e por \cite{soliani2024} (coleta dos dados, pré-processamento, seleção de modelos e avaliação de algoritmos). De forma, que serão adaptadas ao contexto de aprendizado não supervisionado deste projeto.

\subsection{Aquisição de dados}

Serão utilizadas três fontes públicas, todas em nível municipal:

\begin{itemize}
\item \textbf{Secretaria de Segurança Pública do Estado de São Paulo (SSP-SP)}: disponibiliza séries estatísticas mensais, por município, sobre a incidência de diferentes naturezas de ocorrência (e.g., homicídio doloso, roubo, furto, roubo de veículo e estupro), em formato tabular (XLS/CSV).

\item \textbf{Índice de Efetividade da Gestão Municipal (IEGM), do Tribunal de Contas do Estado de São Paulo (TCE-SP)}: fornece, anualmente, o índice geral de efetividade e as notas de suas sete dimensões (educação, saúde, planejamento, gestão fiscal, meio ambiente, cidades e governança em tecnologia da informação).

\item \textbf{Sistema IBGE de Recuperação Automática (SIDRA)}: reúne indicadores socioeconômicos e demográficos consolidados pelo IBGE. Os seguintes indicadores são de interesse para esta análise:  população residente, Produto Interno Bruto (PIB) per capita e taxa de urbanização.

\end{itemize}

\subsection{Pré-processamento}
Após a obtenção dos dados, será realizada a limpeza dos dados (tratamento de valores ausentes, duplicidades e inconsistências de nomenclatura municipal). As ocorrências criminais serão convertidas em taxas por 100 mil habitantes, de modo a permitir a comparação entre municípios de portes populacionais muito distintos, e agregadas em base anual para reduzir a sensibilidade a flutuações mensais de curto prazo, de maneira a seguir a prática adotada em estudos semelhantes de criminalidade municipal no Brasil \cite{soaresfilho2024}. Em seguida, será avaliada a correlação entre as variáveis candidatas.

\subsection{Aplicação e comparação de algoritmos de clusterização}
Serão analisados diferentes algoritmos e comparados sobre uma mesma base de dados pré-processada: \textbf{K-means}, \textbf{agrupamento hierárquico aglomerativo} e \textbf{DBSCAN}. Para o K-means e o agrupamento hierárquico, o número de grupos será determinado a partir do método do cotovelo, do coeficiente de silhueta e do índice de Calinski--Harabasz, testando um intervalo de valores de $k$; para o DBSCAN, os parâmetros de raio de vizinhança e número mínimo de pontos serão definidos por meio da técnica do gráfico de distância ao $k$-ésimo vizinho mais próximo. 

Os agrupamentos obtidos pelos diferentes algoritmos serão comparados entre si por meio das métricas de validação interna (silhueta, Calinski--Harabasz, Davies--Bouldin), além de uma análise qualitativa da estabilidade dos grupos frente a pequenas variações na base de dados. O algoritmo e a configuração que apresentarem o melhor equilíbrio entre qualidade estatística e interpretabilidade prática serão selecionados para a etapa de caracterização dos perfis.

Os experimentos serão implementados em Python, utilizando as bibliotecas \emph{pandas} e \emph{geopandas} para manipulação de dados, \emph{scikit-learn} para os algoritmos de clusterização e métricas de validação, e \emph{matplotlib}/\emph{seaborn} para visualização.


\subsection{Caracterização de perfis}

Os grupos resultantes serão caracterizados por meio de estatísticas descritivas (médias e medianas por variável em cada grupo), gráficos comparativos (\emph{boxplots} e gráficos de radar) e testes estatísticos de comparação entre grupos (ANOVA ou, em caso de não normalidade, Kruskal--Wallis). 

Além disso, os perfis serão representados em mapas coropléticos, para identificar eventuais padrões de contiguidade espacial entre municípios de um mesmo grupo, de forma análoga às representações espaciais utilizadas por Soares Filho et al. \cite{soaresfilho2024} e por Minamisava et al. \cite{minamisava2009}. 


\section{Cronograma}
\begin{table}[ht]
\centering
\caption{Cronograma do trabalho. As linhas em amarelo representam as entregas obrigatórias propostas.}
\label{tab:exTable1}
\includegraphics[width=1.05\textwidth]{Imagens/Cronograma.jpeg}
\end{table}



\bibliographystyle{sbc}
\bibliography{sbc-template}

\end{document}
trabalho